\section{导读：这不是一场模型发布会，而是一张 2026 年 AI 研究的横截面}

这期访谈的主角是姚顺宇：清华物理本科、斯坦福理论物理博士，短暂进入 Berkeley 做博士后后转向 AI，先后在 Anthropic 与 Google DeepMind 做前沿模型训练工作。访谈题目虽然带着人物故事色彩，但真正有价值的部分，是它把 2026 年初模型实验室内部的一组问题串了起来：模型能力是否同质化、预训练是否到头、coding 与 long-horizon agent 为什么突然爆发、蒸馏的边界在哪里、Anthropic 与 Gemini 的组织差异，以及为什么姚顺宇认为 AI 时代已经不再是个人英雄主义的时代。

\begin{importantbox}{本笔记的阅读方式}
本笔记不是逐字转写，而是把 3 小时 48 分钟访谈压缩成一份结构化技术笔记。每一章都保留访谈中的核心判断，并补充必要背景：benchmark、scaling law、agentic coding、distillation、long horizon、TPU/GPU 拓扑、组织设计等。人物经历部分服务于理解研究品味和选择，不会展开成传记。
\end{importantbox}

\begin{figure}[H]
\centering
\includegraphics[width=0.88\textwidth]{frame-two-yaos.jpg}
\caption{开场介绍：两位同名的姚顺宇/姚顺雨，以及本期嘉宾从理论物理转入 AI 的履历。\protect\footnotemark}
\end{figure}
\footnotetext{视频画面时间区间：00:04:12--00:04:40。画面本身是访谈场景，教学价值在于定位人物与问题背景。}

\begin{knowledgebox}{读图：封面与人物帧的作用}
本期不是幻灯片课程，视频画面主要提供语境证据：谁在说、在哪一段说、讨论对象是否已经切换。读者不应从人物画面本身推断技术结论，而应把它当作章节锚点，技术判断来自字幕、章节和上下文整理。
\end{knowledgebox}


\begin{knowledgebox}{访谈主线}
访谈从“两个姚顺宇”的身份辨析开始，迅速进入 AI 行业判断：三大模型实验室在纸面 benchmark 上接近之后，差异转向问题定义、数据构造、产品形态和组织执行。随后，姚顺宇用自己的物理训练经历解释为什么他偏好困难问题、为什么重视实验验证，最后落到组织与集体主义：现代大模型已经太复杂，单个英雄很难独立推动系统级跃迁。
\end{knowledgebox}

\subsection{章节地图}

\noindent\begin{longtable}{p{0.16\textwidth}p{0.32\textwidth}p{0.43\textwidth}}
\toprule
时间段 & 原视频章节 & 本笔记处理重点 \\
\midrule
00:00--00:07 & 两个 Shunyu Yao & 嘉宾履历、物理转 AI、两位同名研究者的差异 \\
00:07--00:25 & 竞争与逃逸 & benchmark 同质化、模型差异从能力转向行为定义 \\
00:25--00:35 & Pre-train 没有到头 & scaling law、数据墙、预训练学习能力 \\
00:35--00:54 & Coding 的爆发 & agentic coding、长程任务、产品 wrapper 与模型壁垒 \\
00:54--01:08 & 蒸馏与机器人 & 硬蒸/软蒸、资源劣势下的中国模型策略、机器人机会 \\
01:08--01:52 & 物理经历 & 非厄米系统、高能物理、物理训练对 AI 实验观的影响 \\
01:52--02:41 & Anthropic & Claude 训练、top-down 机制、产品化与 Claude Code \\
02:41--03:01 & Gemini & ML coding、long horizon、有限训练与近似无限使用 \\
03:01--03:48 & 预测、组织、集体主义 & neo lab、C 端产品、TPU/GPU、AI 研究员素质、个人英雄主义之后 \\
\bottomrule
\end{longtable}

\subsection{本章小结}

本期访谈最重要的信号不是“某家公司更强”，而是前沿模型训练已经进入一个更难外部观察的阶段：公开指标差距缩小，真实差异藏在数据、评估、行为定义、工具使用、长期任务和组织执行里。理解这期访谈，关键是把它当作“AI 工程组织如何在高不确定性下做 bet”的案例。

